\section{Background and Problem Setting}
\label{sec:background}

\paragraph{Lexical statistical watermarks.}
A watermark in the class we study is defined by: (i) a keyed deterministic
\emph{word predicate} $g_K(w) \in \{0,1\}$ (e.g., the top bit of
$\mathrm{HMAC}(K, w)$); (ii) a keyed \emph{band map} $b_K(w) \in [0, B)$
partitioning dictionary words into $B$ frequency bands; (iii) an embedder
that rewrites occurrences of dictionary words toward the target color of
their band (e.g., choosing the synonym mapped to the band's target bit);
(iv) a detector that tokenizes suspect text, counts per-band green votes
$g_b$ and total votes $n_b$, and computes per-band $z_b = (g_b - p_0 n_b) /
\sqrt{p_0(1-p_0)\,n_b}$ with base rate $p_0$ (default 0.5, optionally
calibrated). The \emph{existence score} is $\sum_b |z_b|$; a multi-user
identifier is decoded from band signs and matched against a registry,
with an abstain protocol for low-margin or missing-truth cases.

\paragraph{Threat model.}
The attacker knows the scheme and that a watermark may be present, but not
the key. They may edit leaked text (insert, delete, synonym-swap, reorder)
or \emph{construct} text to be scanned (form letters, template spam). The
verifier holds the key and optional per-document metadata (salt, band set).
We do not consider key compromise; a symmetric key leak is total for any
scheme in this class.

\paragraph{Evaluation protocol and its blind spot.}
Standard protocol: calibrate thresholds on a same-domain null corpus, test
on held-out same-domain samples, report TPR at low FPR. All three of our
findings live \emph{outside} this protocol: repetition attacks
(Finding~1), domain mismatch (Finding~2), and constructed template texts
(Finding~3) are neither in the calibration corpus nor in the test split.
This is why they survive peer review and ship to production.
